\exercisesection{7}

\begin{enumerate}
\def\labelenumi{\arabic{enumi}.}
\item
  设 \((\phi_{t})_{t\in I}\) 为域 \(K\) 的一族位，其值取于有限个域；进一步假设对某个 \(x\in K\)，\(\phi_t(x)\) 所成集合是有限的。证明存在形如第 1 节 (1) 的多项式 \(f(X)\)，使得 \(f(x)\neq 0\)，且元素 \(z = f(x)^{-1}\) 具有如下性质：(i) \(\phi_t(x) = \infty\) 蕴含 \(\phi_t(xz) = 0\) 且 \(\phi_t(z) = 0\)；(ii) \(\phi_t(x)\neq \infty\) 蕴含 \(\phi_t(xz)\neq \infty\) 且 \(\phi_t(z)\neq 0\)。（证明与引理 1 相同。）
\item
  用第 1 节命题 1 的假设与记号，令 q 为 B 的素理想；证明 \(B_{q}\) 是赋值环。
\item
  令 （\(\mathbf{A}_i\)）为域 \(1 \leqslant i \leqslant n\) 的两两独立赋值环，并令 \(\mathbf{K}\)。对每个 \(\mathbf{A} = \bigcap_{i} \mathbf{A}_i\)，令 \(i\) 为 \(a_i\) 的理想 ，并记 \(\neq 0\)；\(\mathbf{A}_i\)\(a = \bigcap_{i} a_i\)；
\end{enumerate}

证明对所有  都有 \(i\)。反之，若 \(\mathfrak{a}_i = \mathbf{A}_i\mathfrak{a}\) 是 \(\mathfrak{b}\) 的理想 ，且 \(\neq 0\)，则 \(\mathbf{A}\)\(\mathfrak{b}_i = \mathbf{A}_i\mathfrak{b}\)（使用第 2 节的定理 1）。\(\mathfrak{b} = \bigcap_i \mathfrak{b}_i\)（使用第 2 节的定理 1）。

